\chapter{Teen aggression}
\index{Teen aggression}

\medskip
\noindent\textit{Educational reference; always seek professional advice for individual care.}

\section*{Definition and Overview}
Aggression in teenagers refers to behaviour that is intended to cause harm, and it can take verbal, physical, or relational forms. Verbal aggression includes yelling, name-calling, and threats; physical aggression includes hitting, pushing, and damaging property; relational aggression includes gossip, exclusion, and damaging friendships. Some anger and conflict are a normal part of adolescence, as young people push for independence while their brains and emotions change rapidly. Aggression becomes a concern when it is frequent, severe, out of proportion, directed at family, peers, or teachers, or when it harms relationships, school life, or safety. Aggression is not simply "bad behaviour": it is almost always a signal of distress, and understanding what sits underneath it is the first step to helping. The causes can include stress, mental-health conditions, neurodevelopmental differences, substance use, and trauma. This chapter, written for parents, carers, teachers, and young people themselves, explains why some teenagers act aggressively, how to tell normal grumpiness from a problem that needs help, what can be done at home and with professionals, and where to turn in a crisis. A calm, non-judgemental response almost always helps more than punishment alone.

\section*{Epidemiology}
Anger and conflict are very common in adolescence. Most parents of a teenager will observe periods of irritability, arguing, and slammed doors, and these peaks are most typical in the early-to-middle teenage years, when the gap between the brain's emotional system and its self-control system is at its widest. More serious aggression is much less common, and physical aggression is more common in boys while relational aggression is more common in girls. A much smaller group, roughly five percent or fewer, shows persistent, serious aggression that meets the criteria for conditions such as conduct disorder. Aggression tends to decline after the teenage years for most young people, but for a minority it becomes a lasting pattern. Rates are higher among teenagers who experience violence at home, who are bullied or excluded, who use alcohol or other drugs heavily, or who have untreated mental-health conditions or neurodevelopmental difficulties such as attention deficit hyperactivity disorder (ADHD). Statistics describe groups, not individuals: a teenager who is aggressive is not doomed to a difficult adulthood, and early, calm help changes the trajectory for the better.

\section*{Aetiology and Pathophysiology}
Teen aggression rarely has a single cause. It arises from a combination of brain development, emotions, circumstances, and sometimes medical or psychiatric conditions.

\begin{itemize}
    \item \textbf{Brain development:} the emotional centres of the teenage brain mature earlier than the prefrontal cortex that governs impulse control, so emotional reactions can outrun reason
    \item \textbf{Hormones:} the surges of puberty can amplify emotional reactivity, though hormones alone do not cause aggression
    \item \textbf{Stress and overwhelm:} school pressure, family conflict, and social demands can overwhelm a teenager's coping system, and aggression becomes a way of discharging distress
    \item \textbf{Trauma and abuse:} teenagers who have experienced violence, abuse, or neglect are more likely to respond to threat with aggression, because their threat system is on alert
    \item \textbf{Mental-health conditions:} depression, anxiety, ADHD, and oppositional defiant or conduct disorders are associated with irritability and aggression, and treating the underlying condition often calms the behaviour
    \item \textbf{Neurodevelopmental differences:} autism and ADHD can make it harder to read social cues, manage frustration, and control impulses
    \item \textbf{Substance use:} alcohol and other drugs lower inhibition and judgement and are strongly linked to aggressive episodes
    \item \textbf{Family and social learning:} teenagers learn from the models around them, and harsh or inconsistent discipline teaches aggressive responses
    \item \textbf{Bullying and exclusion:} being bullied, or struggling with cyberbullying and social rejection, can turn a teenager's pain outward in anger
\end{itemize}

Understanding these drivers matters because the treatment differs: a depressed teenager needs depression treatment, a teenager who has experienced trauma needs safety and trauma-informed care, and an overwhelmed teenager needs support with the load, not just consequences.

\section*{Clinical Features}
Aggression in teenagers shows up in a range of behaviours, and the pattern matters more than any single incident.

\begin{itemize}
    \item \textbf{Verbal aggression:} shouting, swearing, name-calling, and threats directed at parents, siblings, or peers
    \item \textbf{Physical aggression:} hitting, kicking, pushing, throwing or breaking objects, and fighting
    \item \textbf{Relational aggression:} spreading rumours, excluding others, and using friendships as weapons, more common in girls
    \item \textbf{Impulsivity and poor frustration tolerance:} exploding over small frustrations, and seeming unable to pause before acting
    \item \textbf{Irritability and low mood:} frequent grumpiness, a short fuse, and a generally negative mood that colours most interactions
    \item \textbf{Defiance and rule-breaking:} refusing requests, arguing with rules, and breaking curfews or school rules
    \item \textbf{Sleep and appetite changes:} difficulty sleeping, staying up very late, or loss of appetite alongside the behaviour change
    \item \textbf{School problems and withdrawal:} falling grades, detentions, exclusion, and dropping out of activities
    \item \textbf{Self-harm or threats to self:} aggression can turn inward too, and cutting or talk of suicide is a red flag requiring urgent help
\end{itemize}

\section*{Differential Diagnosis}
Because aggression is a symptom, not a diagnosis, professionals look for the condition underneath it.

\begin{itemize}
    \item \textbf{Normal teenage conflict:} arguing and grumpiness that are proportionate, brief, and do not harm relationships are normal; the distinction is frequency, intensity, and impact
    \item \textbf{Oppositional defiant disorder:} a persistent pattern of arguing, defiance, and irritability that is excessive for the age and causes clear problems at home or school
    \item \textbf{Conduct disorder:} more serious, with aggression, rule-breaking, cruelty, and disregard for others, needing structured treatment
    \item \textbf{Depression:} irritability and anger are core symptoms of depression in teenagers, often with low mood, withdrawal, and changes in sleep and appetite
    \item \textbf{Anxiety:} a teenager overwhelmed by worry can lash out when pushed, and anxiety can look like anger
    \item \textbf{ADHD:} impulsivity and emotional dysregulation lead to outbursts that are usually brief, regretted, and linked to difficulty with self-control
    \item \textbf{Autism:} sensory overload and difficulty reading social situations can cause meltdowns that look like aggression but are distress responses
    \item \textbf{Trauma-related reactions:} hypervigilance and quick anger are common after trauma, and the behaviour is a threat response rather than cruelty
    \item \textbf{Substance use:} intoxication and withdrawal both cause mood swings and aggression, and the pattern often follows binges
    \item \textbf{Physical illness:} rarely, thyroid problems, head injury, or neurological conditions cause irritability and aggression, so a medical review has a place
\end{itemize}

\section*{When to Seek Medical Care}
A single aggressive outburst does not usually need professional help, but several situations do. See a GP if aggression is frequent, escalating, or dangerous; if it is accompanied by depression, anxiety, or self-harm; if school or friendships are suffering; if sleep or appetite have changed; or if the family feels they cannot manage. A GP can assess physical causes, screen for mental-health conditions, and refer to a psychologist or child and adolescent psychiatrist. For non-urgent support, parents and young people can call a a local youth crisis service or a a local mental-health support service, and a a local crisis service is available around the clock.

\textbf{Contact your local emergency services for:}
\begin{itemize}
    \item Any immediate risk of serious violence, or a threat to harm another person
    \item Talk of suicide, or self-harm needing medical attention
    \item Behaviour that is dangerous to the teenager themselves, such as running into traffic
    \item Confusion, hallucinations, or a sudden severe change in mental state
\end{itemize}

\section*{Diagnosis and Investigations}
Assessing a teenager with aggression is a careful, non-judgemental process that looks at the whole picture rather than reaching for a label.

\begin{itemize}
    \item \textbf{Full history:} the GP or psychologist asks about the pattern of behaviour, what triggers it, what happens before and after, and how it affects school, friends, and family
    \item \textbf{Development and school review:} how the teenager is functioning academically and socially, and how the school sees the behaviour
    \item \textbf{Family and life circumstances:} recent changes, conflicts, losses, or stressful events at home or among peers
    \item \textbf{Screening questionnaires:} validated tools such as the Strengths and Difficulties Questionnaire help map symptoms, but are only aids to a clinical assessment
    \item \textbf{Medical review:} a physical check and, if indicated, blood tests to exclude thyroid or other physical causes, plus a review of sleep and medicines
    \item \textbf{Substance-use assessment:} a confidential, non-judgemental discussion about alcohol and drug use
    \item \textbf{Specialist referral:} persistent or severe patterns may be referred to a child and adolescent mental-health service or a specialist paediatrician
\end{itemize}

\section*{Management}
Helping a teenager with aggression works best when the whole family is involved, and when the teenager is treated with respect rather than as the "problem".

\begin{itemize}
    \item \textbf{Stay calm and de-escalate:} when a situation is heating up, reducing demands, giving space, and speaking calmly prevents escalation; consequences are applied after, not during, the heat of the moment
    \item \textbf{Set consistent, fair boundaries:} clear expectations about respectful behaviour, with consequences that are predictable, proportionate, and never physical
    \item \textbf{Use psychoeducation:} helping parents and teenagers understand that aggression is often distress or impulse, not malice, changes how everyone responds
    \item \textbf{Cognitive behavioural therapy (CBT):} helps teenagers recognise the thoughts and feelings that lead to anger and learn alternative responses
    \item \textbf{Family therapy:} improving communication and reducing conflict between teenagers and parents addresses the environment in which aggression grows
    \item \textbf{Treat the underlying condition:} if depression, anxiety, ADHD, or another condition is present, treating it is often the most effective way to calm the aggression, and may include medication under specialist guidance
    \item \textbf{Parenting support:} programs such as Triple P and the Incredible Years give parents practical, evidence-based tools for managing challenging behaviour
    \item \textbf{Collaborate with school:} working with the school on supports, behavioural plans, and the pressures the teenager faces at school
    \item \textbf{Safety planning:} if aggression involves violence or risk to others, a structured safety plan is essential, and family violence services should be involved where relevant
\end{itemize}

\section*{Complications}
Unmanaged aggression can damage the teenager's life and the family far beyond the moment.

\begin{itemize}
    \item \textbf{Family conflict and breakdown:} aggression corrodes the trust and warmth families need, and chronic conflict harms everyone
    \item \textbf{School failure and exclusion:} detentions, suspensions, and expulsion set a teenager behind academically and socially
    \item \textbf{Damaged friendships:} aggressive behaviour drives peers away and can lead to the very isolation the teenager is reacting to
    \item \textbf{Legal consequences:} physical violence, vandalism, and threats can lead to police involvement and a criminal record
    \item \textbf{Injury:} fights and angry destruction of property can hurt the teenager and others
    \item \textbf{Substance use:} aggression and substance use often feed each other
    \item \textbf{Long-term conduct problems:} for a minority, persistent aggression predicts antisocial behaviour in adulthood
    \item \textbf{Self-harm and suicide:} aggression turned inward is dangerous, and irritability with self-harm or suicidal thoughts needs urgent professional help
\end{itemize}

\section*{Prognosis and Outcomes}
The outlook for teenage aggression is genuinely good for most young people. A large proportion of aggressive behaviour in adolescence is a phase driven by brain development, hormones, and stress, and it settles as the teenager matures, learns better skills, and gains more control over their life. When the behaviour is a symptom of depression, anxiety, ADHD, or trauma, treating the underlying condition typically calms the aggression as well, and the response to treatment is often rapid and rewarding. The strongest predictor of a good outcome is early, respectful help: teenagers who feel understood respond far better than those who feel punished and lectured. A minority with serious conduct problems benefit from structured, intensive family and behavioural programs, and many still change course. Parents and carers should hold onto the fact that this stage, however hard, is not the finished story of the young person, and that patience, consistency, and professional support turn a difficult phase into a passing one.

\section*{Prevention and Self-Care}
\begin{itemize}
    \item Model calm responses: how you handle your own frustration teaches your teenager how to handle theirs
    \item Protect sleep, regular meals, and exercise, because a tired, hungry teenager is a combustible one
    \item Keep routines and predictable boundaries, and praise cooperative behaviour warmly when it happens
    \item Spend regular, pressure-free time with your teenager so that conflict is not the only way you connect
    \item Teach and name feelings, and practise calming skills such as breathing, walking away, and taking a break
    \item Stay connected to school and teachers, and respond early to signs of bullying or social distress
    \item Look after your own wellbeing, and seek help for yourself if the strain of parenting is overwhelming
\end{itemize}

\section*{Sources and Further Reading}
The following organisations provide reliable information and support for parents and young people dealing with aggression and difficult behaviour.

\begin{itemize}
    \item Raising Children Network. \textit{Teenagers, anger, and aggressive behaviour.} https://raisingchildren.net.au/
    \item Kids Helpline. \textit{Anger and aggressive feelings in teenagers.} https://kidshelpline.com.au/
    \item Beyond Blue. \textit{Teenage mental health and irritability.} https://www.beyondblue.org.au/
    \item Lifeline. \textit{Crisis support for young people.} https://www.lifeline.org.au/
    \item healthdirect Australia. \textit{Anger management and teen behaviour support.} https://www.healthdirect.gov.au/
    \item Headspace. \textit{Mental health support for young Australians.} https://headspace.org.au/
\end{itemize}

\clearpage
